% Options for packages loaded elsewhere
\PassOptionsToPackage{unicode}{hyperref}
\PassOptionsToPackage{hyphens}{url}
\PassOptionsToPackage{table}{xcolor}
%
\documentclass[
  letterpaper,
]{article}
\usepackage{amsmath,amssymb}
\usepackage{iftex}
\usepackage[T1]{fontenc}
\usepackage[utf8]{inputenc}
\usepackage{textcomp}
\usepackage{lmodern}
\ifPDFTeX\else
  % xetex/luatex font selection
\fi
% Use upquote if available, for straight quotes in verbatim environments
\IfFileExists{upquote.sty}{\usepackage{upquote}}{}
\IfFileExists{microtype.sty}{% use microtype if available
  \usepackage[]{microtype}
  \UseMicrotypeSet[protrusion]{basicmath} % disable protrusion for tt fonts
}{}
\makeatletter
\@ifundefined{KOMAClassName}{% if non-KOMA class
  \IfFileExists{parskip.sty}{%
    \usepackage{parskip}
  }{% else
    \setlength{\parindent}{0pt}
    \setlength{\parskip}{6pt plus 2pt minus 1pt}}
}{% if KOMA class
  \KOMAoptions{parskip=half}}
\makeatother
\usepackage{xcolor}
\usepackage{color}
\usepackage{fancyvrb}
\newcommand{\VerbBar}{|}
\newcommand{\VERB}{\Verb[commandchars=\\\{\}]}
\DefineVerbatimEnvironment{Highlighting}{Verbatim}{commandchars=\\\{\}}
% Add ',fontsize=\small' for more characters per line
\newenvironment{Shaded}{}{}
\newcommand{\AlertTok}[1]{\textcolor[rgb]{1.00,0.00,0.00}{\textbf{#1}}}
\newcommand{\AnnotationTok}[1]{\textcolor[rgb]{0.38,0.63,0.69}{\textbf{\textit{#1}}}}
\newcommand{\AttributeTok}[1]{\textcolor[rgb]{0.49,0.56,0.16}{#1}}
\newcommand{\BaseNTok}[1]{\textcolor[rgb]{0.25,0.63,0.44}{#1}}
\newcommand{\BuiltInTok}[1]{\textcolor[rgb]{0.00,0.50,0.00}{#1}}
\newcommand{\CharTok}[1]{\textcolor[rgb]{0.25,0.44,0.63}{#1}}
\newcommand{\CommentTok}[1]{\textcolor[rgb]{0.38,0.63,0.69}{\textit{#1}}}
\newcommand{\CommentVarTok}[1]{\textcolor[rgb]{0.38,0.63,0.69}{\textbf{\textit{#1}}}}
\newcommand{\ConstantTok}[1]{\textcolor[rgb]{0.53,0.00,0.00}{#1}}
\newcommand{\ControlFlowTok}[1]{\textcolor[rgb]{0.00,0.44,0.13}{\textbf{#1}}}
\newcommand{\DataTypeTok}[1]{\textcolor[rgb]{0.56,0.13,0.00}{#1}}
\newcommand{\DecValTok}[1]{\textcolor[rgb]{0.25,0.63,0.44}{#1}}
\newcommand{\DocumentationTok}[1]{\textcolor[rgb]{0.73,0.13,0.13}{\textit{#1}}}
\newcommand{\ErrorTok}[1]{\textcolor[rgb]{1.00,0.00,0.00}{\textbf{#1}}}
\newcommand{\ExtensionTok}[1]{#1}
\newcommand{\FloatTok}[1]{\textcolor[rgb]{0.25,0.63,0.44}{#1}}
\newcommand{\FunctionTok}[1]{\textcolor[rgb]{0.02,0.16,0.49}{#1}}
\newcommand{\ImportTok}[1]{\textcolor[rgb]{0.00,0.50,0.00}{\textbf{#1}}}
\newcommand{\InformationTok}[1]{\textcolor[rgb]{0.38,0.63,0.69}{\textbf{\textit{#1}}}}
\newcommand{\KeywordTok}[1]{\textcolor[rgb]{0.00,0.44,0.13}{\textbf{#1}}}
\newcommand{\NormalTok}[1]{#1}
\newcommand{\OperatorTok}[1]{\textcolor[rgb]{0.40,0.40,0.40}{#1}}
\newcommand{\OtherTok}[1]{\textcolor[rgb]{0.00,0.44,0.13}{#1}}
\newcommand{\PreprocessorTok}[1]{\textcolor[rgb]{0.74,0.48,0.00}{#1}}
\newcommand{\RegionMarkerTok}[1]{#1}
\newcommand{\SpecialCharTok}[1]{\textcolor[rgb]{0.25,0.44,0.63}{#1}}
\newcommand{\SpecialStringTok}[1]{\textcolor[rgb]{0.73,0.40,0.53}{#1}}
\newcommand{\StringTok}[1]{\textcolor[rgb]{0.25,0.44,0.63}{#1}}
\newcommand{\VariableTok}[1]{\textcolor[rgb]{0.10,0.09,0.49}{#1}}
\newcommand{\VerbatimStringTok}[1]{\textcolor[rgb]{0.25,0.44,0.63}{#1}}
\newcommand{\WarningTok}[1]{\textcolor[rgb]{0.38,0.63,0.69}{\textbf{\textit{#1}}}}
\usepackage{longtable,booktabs,array}
\usepackage{calc} % for calculating minipage widths
% Correct order of tables after \paragraph or \subparagraph
\usepackage{etoolbox}
\makeatletter
\patchcmd\longtable{\par}{\if@noskipsec\mbox{}\fi\par}{}{}
\makeatother
% Allow footnotes in longtable head/foot
\IfFileExists{footnotehyper.sty}{\usepackage{footnotehyper}}{\usepackage{footnote}}
\makesavenoteenv{longtable}
\setlength{\emergencystretch}{3em} % prevent overfull lines
\providecommand{\tightlist}{%
  \setlength{\itemsep}{0pt}\setlength{\parskip}{0pt}}
\setcounter{secnumdepth}{5}
% Self-contained presentation rules embedded in every generated LaTeX file.
\usepackage{iftex}
\ifPDFTeX\else\PackageError{SSE-ASAS}{Compile with pdfLaTeX only}{}\fi
\usepackage[T1]{fontenc}
\usepackage[utf8]{inputenc}
\usepackage{lmodern}
\usepackage{microtype}
\usepackage{geometry}
\geometry{margin=0.93in,headheight=16pt}
\usepackage{xcolor} % [table] passed at top of file (option-clash fix)
\definecolor{ink}{HTML}{15283B}
\definecolor{brand}{HTML}{075B75}
\definecolor{brandlight}{HTML}{EAF5F8}
\definecolor{warm}{HTML}{A15C1D}
\definecolor{muted}{HTML}{536577}
\usepackage{fancyhdr}
\usepackage{titlesec}
\usepackage{enumitem}
\usepackage{booktabs}
\usepackage{longtable}
\usepackage{array}
\usepackage{tabularx}
\usepackage{tcolorbox}
\tcbuselibrary{breakable,skins}
\usepackage{listings}
\usepackage{fancyvrb}
\usepackage{hyperref}
\hypersetup{colorlinks=true,linkcolor=brand,urlcolor=brand,citecolor=brand}
\urlstyle{same}
\usepackage{lastpage}
\setlength{\emergencystretch}{3em}
\setlength{\parskip}{0.45em}
\setlength{\parindent}{0pt}
\setlist[itemize]{leftmargin=*,topsep=3pt,itemsep=2pt}
\setlist[enumerate]{leftmargin=*,topsep=3pt,itemsep=2pt}
\titleformat{\section}{\Large\bfseries\color{ink}}{\thesection}{0.75em}{}[\vspace{-2pt}\color{brand}\titlerule]
\titleformat{\subsection}{\large\bfseries\color{brand}}{\thesubsection}{0.7em}{}
\titleformat{\subsubsection}{\normalsize\bfseries\color{ink}}{\thesubsubsection}{0.7em}{}
\pagestyle{fancy}
\fancyhf{}
\fancyhead[L]{\small\sffamily\textcolor{brand}{SOFTWARE SECURITY ENGINEERING}}
\fancyhead[R]{\small\sffamily\textcolor{muted}{TECHNICAL SERIES}}
\fancyfoot[L]{\scriptsize\sffamily Version 1.0 \textbullet\ 14 October 2026}
\fancyfoot[R]{\scriptsize\sffamily Page \thepage\ of \pageref{LastPage}}
\renewcommand{\headrulewidth}{0.5pt}
\renewcommand{\footrulewidth}{0.2pt}
\setcounter{tocdepth}{2}
\setcounter{secnumdepth}{3}
\lstdefinestyle{sse}{basicstyle=\ttfamily\footnotesize,frame=single,framerule=0.4pt,rulecolor=\color{brand!40},backgroundcolor=\color{brandlight!45},numbers=left,numberstyle=\scriptsize\color{muted},numbersep=6pt,breaklines=true,breakatwhitespace=false,tabsize=2,showstringspaces=false,columns=fullflexible,keepspaces=true,captionpos=b,aboveskip=9pt,belowskip=9pt}
\lstset{style=sse}
\newtcolorbox{sourceconcept}[1][]{enhanced,breakable,colback=brandlight,colframe=brand!65!black,title={ORIGINAL CONCEPT},fonttitle=\bfseries\sffamily,#1}
\newtcolorbox{modernextension}[1][]{enhanced,breakable,colback=white,colframe=warm!85!black,title={MODERN EXTENSION},fonttitle=\bfseries\sffamily,#1}
\newtcolorbox{fieldwork}[1][]{enhanced,breakable,colback=white,colframe=ink!55,title={PRACTICAL APPLICATION},fonttitle=\bfseries\sffamily,#1}
\AtBeginDocument{\color{ink}}
\makeatletter
\renewcommand{\maketitle}{%
 \begin{titlepage}\thispagestyle{empty}
 \vspace*{1.5cm}{\color{brand}\rule{\textwidth}{3pt}}\par\vspace{1cm}
 {\sffamily\bfseries\color{muted}\large SOFTWARE SECURITY ENGINEERING\\ APPLICATION SECURITY ASSESSMENT SERIES}\par\vspace{1.05cm}
 {\sffamily\bfseries\color{ink}\fontsize{27}{34}\selectfont\@title\par}\vspace{0.8cm}
 {\sffamily\large\color{brand}\@author\par}\vfill
 \begin{tcolorbox}[colback=brandlight,colframe=brand!55,boxrule=0.5pt]
 \small\textbf{Editorial policy}\quad Original teaching synthesis with book-grounded concepts distinguished from contemporary practices. Modern platform behavior is version-sensitive; validate against linked vendor documentation before production use.
 \end{tcolorbox}
 {\color{brand}\rule{\textwidth}{2pt}}\par\vspace{0.3cm}
 {\sffamily\small Version 1.0 \quad|\quad October 2026 \quad|\quad pdfLaTeX}\end{titlepage}\setcounter{page}{1}}
\makeatother
\ifLuaTeX
  \usepackage{selnolig}  % disable illegal ligatures
\fi
\usepackage{bookmark}
\IfFileExists{xurl.sty}{\usepackage{xurl}}{} % add URL line breaks if available
\urlstyle{same}
\hypersetup{
  pdftitle={Software Vulnerability Taxonomy},
  pdfauthor={Independent Technical Study Edition},
  hidelinks,
  pdfcreator={LaTeX via pandoc}}

\title{Software Vulnerability Taxonomy}
\author{Independent Technical Study Edition}
\date{}

\begin{document}
\maketitle

{
\setcounter{tocdepth}{3}
\tableofcontents
}
\section{Why taxonomy matters}\label{why-taxonomy-matters}

A useful vulnerability classification records the \emph{failure
mechanism}, \emph{violated invariant}, \emph{affected trust boundary},
and \emph{verification evidence}. A label such as ``injection'' alone is
insufficient for an engineer deciding whether a finding is real and
whether a fix prevents recurrence.

\textbf{Primary lineage:} Dowd, McDonald, and Schuh, \emph{The Art of
Software Security Assessment}, Chapters 1--3 and 5--18; Chess and West,
\emph{Secure Programming with Static Analysis}, Chapters 1, 5--12;
Bishop, Chapter 24. \textbf{Modern extension:} CWE and OWASP are
contemporary cross-references. A CWE ID is a weakness category, an OWASP
Top 10 entry is a risk family, and a CVE identifies a publicly disclosed
vulnerability instance when one has been assigned.

\section{Two complementary classification
axes}\label{two-complementary-classification-axes}

\subsection{Lifecycle / origin of
failure}\label{lifecycle-origin-of-failure}

\begin{itemize}
\tightlist
\item
  \textbf{Design flaw:} authorization omission, misplaced trust, poor
  isolation, non-atomic approval, insecure protocol assumptions.
\item
  \textbf{Implementation flaw:} integer overflow, unsafe buffer write,
  incorrect escaping, SQL construction, path traversal, unsafe
  deserialization.
\item
  \textbf{Operational weakness:} insecure default, excessive deployment
  privileges, missing rotation, vulnerable runtime configuration, unsafe
  diagnostic endpoint.
\end{itemize}

In practice these categories overlap. An authorization flaw may reflect
both a faulty design and an implementation that omitted an object check.
Record the strongest causal explanation and the relevant contributing
causes rather than forcing mutually exclusive labels.

\subsection{Security property
violated}\label{security-property-violated}

\begin{longtable}[]{@{}
  >{\raggedright\arraybackslash}p{(\columnwidth - 4\tabcolsep) * \real{0.3333}}
  >{\raggedright\arraybackslash}p{(\columnwidth - 4\tabcolsep) * \real{0.3333}}
  >{\raggedright\arraybackslash}p{(\columnwidth - 4\tabcolsep) * \real{0.3333}}@{}}
\toprule\noalign{}
\begin{minipage}[b]{\linewidth}\raggedright
Mechanism
\end{minipage} & \begin{minipage}[b]{\linewidth}\raggedright
Likely property
\end{minipage} & \begin{minipage}[b]{\linewidth}\raggedright
Characteristic trust-boundary failure
\end{minipage} \\
\midrule\noalign{}
\endhead
\bottomrule\noalign{}
\endlastfoot
Broken object authorization & Confidentiality/integrity & User-selected
ID bypasses ownership checks \\
SQL injection & Confidentiality/integrity & Data becomes SQL syntax \\
Command injection & Integrity/availability & Text becomes process
command syntax \\
Cross-site scripting & Confidentiality/integrity & Untrusted data enters
browser-executable context \\
Buffer overwrite & Integrity/availability & Unchecked size exceeds
memory object bounds \\
Race condition & Integrity & Policy decision separated from mutable
state \\
Exposed secrets & Confidentiality & Credential enters unauthorized
distribution channel \\
Deserialization abuse & Integrity & Untrusted data acquires behavior or
object authority \\
SSRF & Confidentiality/integrity & User-supplied target reaches
privileged network context \\
\end{longtable}

\section{Engineering catalog: root cause, evidence,
control}\label{engineering-catalog-root-cause-evidence-control}

\begin{longtable}[]{@{}
  >{\raggedright\arraybackslash}p{(\columnwidth - 4\tabcolsep) * \real{0.3333}}
  >{\raggedright\arraybackslash}p{(\columnwidth - 4\tabcolsep) * \real{0.3333}}
  >{\raggedright\arraybackslash}p{(\columnwidth - 4\tabcolsep) * \real{0.3333}}@{}}
\toprule\noalign{}
\begin{minipage}[b]{\linewidth}\raggedright
Weakness / example CWE
\end{minipage} & \begin{minipage}[b]{\linewidth}\raggedright
Evidence to seek
\end{minipage} & \begin{minipage}[b]{\linewidth}\raggedright
Primary prevention
\end{minipage} \\
\midrule\noalign{}
\endhead
\bottomrule\noalign{}
\endlastfoot
Improper authorization (CWE-862) & Missing server-side action guard &
Centralized authorization plus negative tests \\
IDOR / missing object authorization (CWE-639) & Cross-owner record
returns & Scope reads/writes to current principal \\
SQL injection (CWE-89) & Input influences query grammar & Bound
parameters; restricted DB rights \\
OS command injection (CWE-78) & Data reaches shell parser & Avoid shell;
argument-vector APIs \\
XSS (CWE-79) & Value reaches HTML/JS sink & Contextual output encoding;
safe DOM APIs \\
Path traversal (CWE-22) & Attacker path escapes root & Resolved-path
confinement and safe file APIs \\
Out-of-bounds write (CWE-787) & Write index \textgreater{} allocated
extent & Bounds-safe APIs; checked size arithmetic \\
Integer overflow (CWE-190) & Size computation wraps & Guard arithmetic
before allocation or use checked API \\
TOCTOU race (CWE-367) & Check and use separate on mutable target &
Atomic operation, file descriptor, or lock \\
Uncontrolled resource consumption (CWE-400) & Input amplifies CPU/memory
& Limits, timeouts, quotas, admission control \\
Improper restriction of XXE (CWE-611) & XML parser resolves external
entity & Harden parser; disable external entity processing \\
Hard-coded credential (CWE-798) & Secret committed into artifact &
Managed identity/secret store, rotation \\
Missing authentication (CWE-306) & Sensitive function reachable
unauthenticated & Authn at trustworthy enforcement point \\
Sensitive data exposure through logs (CWE-532) & Token or PII emitted to
logs & Data minimization and log sanitization \\
\end{longtable}

\emph{Mapping note:} CWE identifiers are illustrative, not a claim that
the listed IDs exhaust a weakness. Select the most precise applicable
CWE from its official definition and verify version-dependent names and
abstraction levels.

\section{Deep mechanisms}\label{deep-mechanisms}

\subsection{Integer and memory safety}\label{integer-and-memory-safety}

A size expression \texttt{count\ *\ width} can wrap before allocation.
If a later loop uses the original count, the allocation and write bounds
disagree. In C, signed overflow is undefined behavior; unsigned
arithmetic wraps modulo its width. Detecting a ``short buffer''
therefore requires tracing numerical ranges, conversion rules, and every
bound in the producer/consumer chain. Historic examples from Dowd
Chapters 5--8 and Chess Chapters 6--7 remain valuable because the
mechanisms persist even when mitigation technology changes.

\begin{Shaded}
\begin{Highlighting}[]
\PreprocessorTok{\#include }\ImportTok{\textless{}stdint.h\textgreater{}}
\PreprocessorTok{\#include }\ImportTok{\textless{}stdlib.h\textgreater{}}
\DataTypeTok{int}\NormalTok{ reserve}\OperatorTok{(}\DataTypeTok{size\_t}\NormalTok{ count}\OperatorTok{,} \DataTypeTok{size\_t}\NormalTok{ width}\OperatorTok{,} \DataTypeTok{void} \OperatorTok{**}\NormalTok{out}\OperatorTok{)} \OperatorTok{\{}
    \ControlFlowTok{if} \OperatorTok{(!}\NormalTok{out }\OperatorTok{||} \OperatorTok{!}\NormalTok{width }\OperatorTok{||}\NormalTok{ count }\OperatorTok{\textgreater{}}\NormalTok{ SIZE\_MAX }\OperatorTok{/}\NormalTok{ width}\OperatorTok{)} \ControlFlowTok{return} \OperatorTok{{-}}\DecValTok{1}\OperatorTok{;}
    \OperatorTok{*}\NormalTok{out }\OperatorTok{=}\NormalTok{ malloc}\OperatorTok{(}\NormalTok{count }\OperatorTok{*}\NormalTok{ width}\OperatorTok{);}
    \ControlFlowTok{return} \OperatorTok{*}\NormalTok{out }\OperatorTok{?} \DecValTok{0} \OperatorTok{:} \OperatorTok{{-}}\DecValTok{1}\OperatorTok{;}
\OperatorTok{\}}
\end{Highlighting}
\end{Shaded}

\textbf{Assessment:} arithmetic guard is necessary but not sufficient;
check downstream pointer bounds, initialization, ownership, and
allocation-failure behavior. An allocator failing safely does not prove
later writes are bounded.

\subsection{Trust, interpretation, and
context}\label{trust-interpretation-and-context}

An input string does not need to be ``malicious'' until some interpreter
grants it meaning: SQL parser, shell, HTML template, URL resolver, file
path resolver, or log parser. Trace the transition where data acquires
syntax or authority. Sanitization is context-specific; HTML escaping is
not equivalent to SQL parameterization, and path normalization is not an
authorization decision.

\begin{Shaded}
\begin{Highlighting}[]
\CommentTok{\# Vulnerable: user\_id becomes SQL syntax.}
\NormalTok{sql }\OperatorTok{=} \StringTok{"SELECT * FROM invoices WHERE owner\_id = \textquotesingle{}"} \OperatorTok{+}\NormalTok{ user\_id }\OperatorTok{+} \StringTok{"\textquotesingle{}"}
\CommentTok{\# Correct pattern: value is bound data (driver placeholder syntax varies).}
\NormalTok{cur.execute(}\StringTok{"SELECT * FROM invoices WHERE owner\_id = ?"}\NormalTok{, (user\_id,))}
\end{Highlighting}
\end{Shaded}

\subsection{Time, state, and environmental
assumptions}\label{time-state-and-environmental-assumptions}

TOCTOU occurs when a permission, path, or object state is checked and
then used after it can change. Protect invariant-carrying operations
with atomic facilities, transactions, locks, or descriptor-based APIs.
For business logic, guard state changes inside the authoritative data
store, not solely in a preliminary HTTP check.

\section{Taxonomy-to-framework
crosswalk}\label{taxonomy-to-framework-crosswalk}

The table illustrates \emph{families of concern}, not equivalence
between book terminology, CWE definitions, or OWASP Top 10:2025
categories. Exact OWASP category names and member CWEs must be looked up
in the versioned official release before reporting a regulatory mapping.

\begin{longtable}[]{@{}
  >{\raggedright\arraybackslash}p{(\columnwidth - 4\tabcolsep) * \real{0.3333}}
  >{\raggedright\arraybackslash}p{(\columnwidth - 4\tabcolsep) * \real{0.3333}}
  >{\raggedright\arraybackslash}p{(\columnwidth - 4\tabcolsep) * \real{0.3333}}@{}}
\toprule\noalign{}
\begin{minipage}[b]{\linewidth}\raggedright
Source-domain concept
\end{minipage} & \begin{minipage}[b]{\linewidth}\raggedright
CWE family example
\end{minipage} & \begin{minipage}[b]{\linewidth}\raggedright
Contemporary verification area
\end{minipage} \\
\midrule\noalign{}
\endhead
\bottomrule\noalign{}
\endlastfoot
Trust and design-review failures & CWE-862, CWE-639 & ASVS authorization
requirements \\
Metacharacter interpretation & CWE-89, CWE-78, CWE-79 & ASVS
input/encoding and output contexts \\
Native memory corruption & CWE-787, CWE-190 & Memory-safe implementation
and compiler tooling \\
Privilege and file-system mistakes & CWE-269, CWE-22, CWE-367 & OS and
service authorization controls \\
Secrets/privacy & CWE-798, CWE-532 & Secret scanning, credential
lifecycle \\
Operational insecure deployment & CWE-16 (broad category) &
Configuration and pipeline review \\
\end{longtable}

\begin{modernextension}
The OWASP Top 10:2025 is an awareness and prioritization tool, not a complete test plan. Use ASVS 5.0.0 for verifiable controls, CWE for precise failure mechanisms, and organization-specific impact analysis for severity. CodeQL/other SAST results help find implementations; absence of an alert is not proof against business-logic flaws.
\end{modernextension}

\section{Classification decision
procedure}\label{classification-decision-procedure}

\begin{enumerate}
\def\labelenumi{\arabic{enumi}.}
\tightlist
\item
  Identify the attacker-controlled boundary and the sensitive object or
  operation.
\item
  State expected invariant; demonstrate the violation using authorized
  test conditions.
\item
  Distinguish trigger, contributing assumptions, root cause, and
  outcome.
\item
  Classify design/implementation/operational origin and affected
  security property.
\item
  Map to the most specific appropriate CWE after reviewing its
  definition.
\item
  Determine exposure, privileges, reachable execution path, existing
  mitigations, and consequence.
\item
  Record scanner confidence separately from severity and likelihood of
  exploitation.
\item
  Specify fix and \emph{negative regression} test that would fail before
  and pass after remediation.
\end{enumerate}

\textbf{Example:} A tenant can see another tenant's order by changing
\texttt{/orders/123}. Root cause: absent object-level authorization.
Type: design and/or implementation flaw; property: confidentiality.
CWE-639 may apply when an externally supplied key selects a record
outside ownership scope. Fix: query scoped by tenant, plus server-side
decision. Regression: tenant A receives authorization failure for tenant
B's order even when the identifier is valid.

\section{Verification and reporting
artifacts}\label{verification-and-reporting-artifacts}

For each finding, collect: title; affected component and version;
violated requirement; trace from entry point to sensitive operation;
reproduction in a safe test context; impact; CWE mapping; uncertainty;
remediation; regression test; evidence links; owner and status. Avoid
reporting exploitability as proven merely because a scanner assigned
high severity; avoid treating an unproven positive as a false positive
without inspecting the actual execution path.

\begin{fieldwork}
**Catalog assignment:** select ten repository alerts across at least four weaknesses. Group by root-cause mechanism, assign validated CWE IDs, distinguish confidence and business impact, and produce a single prevention pattern per class with one regression test for each pattern.
\end{fieldwork}

\section{Annex A: detailed failure-pattern
profiles}\label{annex-a-detailed-failure-pattern-profiles}

\subsection{Authorization and workflow
flaws}\label{authorization-and-workflow-flaws}

\textbf{Case T01 -- Confused deputy.} A privileged service performs an
operation for an untrusted caller using the service's ambient authority.
Review whether the caller's scope is carried into the decision or
silently replaced by service-wide privileges. Remediate by binding
delegated authority to actor, resource, action, and audience; add tests
with lower-privilege callers.

\textbf{Case T02 -- Mass assignment.} A framework binds arbitrary
request fields into a persistence model and unintentionally accepts
\texttt{role}, \texttt{status}, or \texttt{tenant\_id}. Enforce field
allowlists in input DTOs and protect authorization-sensitive state
transitions at the domain layer. Verify that forbidden fields are not
changed even when supplied with syntactically valid values.

\textbf{Case T03 -- Business process sequence bypass.} A refund
operation accepts a completed invoice without a required preceding
approval. Document the intended state machine, test invalid transitions,
and enforce conditions atomically at the authoritative state store.

\subsection{Input, interpretation, and encoding
flaws}\label{input-interpretation-and-encoding-flaws}

\textbf{Case T04 -- Ambiguous canonicalization.} One layer validates an
encoded path or URL while a downstream layer decodes it differently.
Trace exact byte and character representations through proxy, framework,
URI parser, filesystem, and storage APIs. Normalize once under a
documented policy and check the final resolved object in the correct
trust domain.

\textbf{Case T05 -- Template injection.} User-controlled content enters
a template expression interpreter rather than a literal data
placeholder. Data escaping alone may not protect code-execution features
of the template engine. Separate templates from data, sandbox where
appropriate, and restrict attacker influence over template source.

\textbf{Case T06 -- SSRF / request routing.} A user-supplied URL can
reach loopback, private service networks, cloud metadata services, or
credentials. Model DNS resolution, redirects, URL parsing differences,
egress firewall policy, and protocol handlers. Host allowlists and
infrastructure-level egress control are complementary.

\textbf{Case T07 -- Log injection and disclosure.} User input can forge
lines or structured events and credentials are serialized into
telemetry. Use structured logging with proper encoding, redact sensitive
fields, restrict log access, and test that malicious newlines cannot
impersonate authoritative events.

\subsection{Native and state-management
flaws}\label{native-and-state-management-flaws}

\textbf{Case T08 -- Use-after-free.} Multiple aliases outlive an
object's allocation and a later access observes reclaimed memory.
Analyze ownership/lifetime and asynchronous callbacks. Prefer
memory-safe abstractions and use sanitizers during testing; a null
assignment to one pointer does not repair all aliases.

\textbf{Case T09 -- Integer truncation.} A 64-bit length is cast into a
smaller buffer size, but a later copy uses the original length. Prove a
checked upper bound before conversion, propagate the bounded type, and
test maximum representable values.

\textbf{Case T10 -- Double-free.} Separate error-handling paths
deallocate the same object. Centralize ownership and cleanup
responsibility, and test multiple failure branches. A heap hardening
layer may abort exploitation without making the defect correct.

\textbf{Case T11 -- TOCTOU filesystem race.} A pre-check uses a pathname
and then a later operation resolves that pathname again. Prefer atomic
opens and descriptor-based validation. An attacker-controllable
directory invalidates naive check-then-open defenses.

\textbf{Case T12 -- Non-idempotent asynchronous operation.} Re-delivered
webhook triggers duplicate financial action. Enforce idempotency keys,
unique constraints, and transaction boundaries; asynchronous retries
must not silently bypass business invariants.

\subsection{Configuration and supply-chain
flaws}\label{configuration-and-supply-chain-flaws}

\textbf{Case T13 -- Privileged build automation.} A CI job executes
untrusted pull-request code while holding high-privilege repository
tokens. Separate trusted and untrusted events, minimize permissions, pin
trusted dependencies as policy requires, and constrain self-hosted
runners and cache sharing.

\textbf{Case T14 -- Excessive runtime permission.} An application only
needs object-read access but its service credential can alter policy or
enumerate all customer objects. Review effective rather than nominal
permissions and bind identities to workload and resource.

\textbf{Case T15 -- Missing inventory of code not scanned.} CodeQL
reports zero alerts while the build extracted no relevant language, or a
path filter skipped the package. This is primarily a
\emph{coverage-control failure}, not a software vulnerability in that
package. Record tooling gaps separately and require a successful rerun
before asserting any analyzed scope.

\subsection{Structured classification
exercise}\label{structured-classification-exercise}

For any three cases above, write a complete five-line decision: security
invariant; attacker preconditions; source/operation trace; likely CWE to
validate on the official MITRE site; and the smallest regression test.
Document whether the case arises at design, implementation, operation,
or multiple layers. Never infer severity solely from a named weakness: a
vulnerability's reachability and business consequences can differ
materially even within one CWE class.

\section{Source and modern extension
map}\label{source-and-modern-extension-map}

\textbf{Book-grounded:} Dowd et al., Chapters 1--4 (vulnerability
fundamentals, design, operations, review), 5--8 (memory, C, strings),
9--13 (platform/concurrency), 14--18 (protocols and applications); Chess
\& West, Chapters 1--4, 5--12; Bishop, Chapter 24.

\textbf{Modern extension:} CWE \url{https://cwe.mitre.org/}; OWASP Top
10:2025 \url{https://top10.owasp.org/2025/}; OWASP ASVS 5.0.0
\url{https://owasp.org/projects/asvs}; CodeQL query catalog
\url{https://docs.github.com/en/code-security/reference/code-scanning/codeql/codeql-queries}.

\end{document}
